\section{Introduction}
\label{sec:intro}

Hourly Ethereum (ETH) returns carry a marked downside tail asymmetry. Over
2021--2025 the skewness is \SkewEth{} with excess kurtosis of \KurtEth{}, the
worst hour ($\MinRet\%$) is far larger than the best ($+\MaxRet\%$), and hours
below $-5\%$ outnumber hours above $+5\%$ two to one
(\FarFiveDown{} versus \FarFiveUp).\footnote{The asymmetry is a
\emph{far-extreme} property: counts beyond absolute $\pm 5\%$/$\pm 7\%$
thresholds. The 1\% and 5\% quantiles themselves are nearly symmetric (ratio
\QratioOnePct{} at the 1\% level), a distinction that matters below. The
sign of measured crypto skewness differs across return definitions and
samples: in simple returns it is positive across major cryptocurrencies
\citep{liu2021}, while the hourly log returns studied here skew negative. Fat tails are
the expected stylized fact \citep{cont2001,osterrieder2017}.}
For equities, the textbook reading of such asymmetry runs through the
leverage effect. In crypto, that channel is famously absent or inverted
\citep{baur2018,brini2022,aharon2023}, and the literature falls back on a
different mechanism: the liquidation cascade of the leverage cycle
\citep{geanakoplos2010,brunnermeier2009}, in which falling prices force
collateralized positions to unwind, and forced selling begets further falls.
The mechanism is plausible and widely invoked, both in empirical studies of
DeFi leverage and liquidation risk \citep{perez2021} and in official-sector
financial-stability assessments \citep{oecd2023,fsb2023}. In centralized venues it is essentially
unobservable.

Decentralized finance changes the observability. DeFi lending protocols
liquidate undercollateralized positions mechanically, on-chain, transaction
by transaction. The shadow quantity at the heart of equilibrium leverage
models, the tightening of collateral constraints, is recorded in public
data. We assemble an hourly panel of ETH-collateralized DeFi liquidations
(\LiqTotalBn{} billion USD across 2021--2025) and ask whether this observable
forced deleveraging amplifies the downside tail of ETH returns, as the
leverage-cycle prior predicts. We are explicit about the design. The
measure is aggregated to the hour, covers DeFi only (centralized-exchange
liquidations are out of scope), and at short horizons the liquidation--return
link is mechanically bidirectional: \PLiqWorstOne\% of the worst-1\% return
hours contain contemporaneous liquidations, against \PLiqBestOne\% of the
best-1\% hours and \PLiqAll\% unconditionally. The register of the paper is
therefore descriptive and predictive, never causal.

Our empirical vehicle is the quantile local projection (QLP) of
\citet{han2024}, the natural transposition of the Growth-at-Risk logic of
\citet{adrian2019}, in which tails can respond where the mean is silent, to
high-frequency crypto returns. The naive QLP delivers what the prior
predicts, and what we suspect this nascent literature will soon be reporting.
On impact, the response of the first return percentile to a liquidation shock
is \RatioTailMedImpact$\times$ the (near-zero) median response; the
coefficient deepens monotonically to \BetaQoneHtwelve{} at twelve hours; the
left tail appears to respond more than the right; and the pattern survives a
robustness battery with Bonferroni adjustments. It looks like
textbook endogenous fragility. It is not. Showing why, with a
diagnostic battery that travels beyond crypto, is the contribution of this
paper.

The diagnostics dismantle the downside-amplification reading and replace it
with something more modest and, we argue, more useful. First, a
\emph{symmetric} placebo, a sign-flip world that preserves the exact
empirical volatility path and the full magnitude distribution of returns but
imposes zero skew by construction \citep{wu1986,liu1988,davidson2008},
reproduces the entire left--right tail gap at every horizon, including the
deepening profile. The apparent downside specificity is what symmetric
heteroskedasticity looks like through the lens of extreme quantiles, the
mechanism identified in macroeconomic data by \citet{carriero2024} and, for
scale versus shape in quantile regression, by \citet{koenker1982} and
\citet{machado2000}. Second, the clean per-period counterpart of the naive
regression, tail-exceedance indicators in the spirit of
\citet{engle2004}, shows the shock loads \emph{equally on both tails}
($\Delta = \beta^{\text{down}} - \beta^{\text{up}} = \DeltaExcImpact$ at the
1\% level on impact, $p = \DeltaExcImpactP$), while the overlapping cumulative
variant manufactures a significant asymmetry at twelve hours: the artifact
made visible in a single comparison \citep{hodrick1992,boudoukh2008,
neuberger2021,farago2023}. Third, re-running the full QLP with
Bitcoin as the outcome reproduces the ETH signature at
$\BtcOverEthHtwelve{}$ of its magnitude at twelve hours, with the same
deepening: the pattern is generic crypto stress, not an ETH-specific DeFi
channel. Fourth, an equivalence test in the spirit of
\citet{altman1995} and \citet{lakens2017} converts the null into a bound.
Net of volatility, we can exclude downside-specific amplification
larger than \MdeEightyExc{} per unit of the log-liquidation shock, below
our pre-specified smallest effect size of interest under all three
calibrations of the shock span, including the strictest. This bound is a
one-percent-tail statement; it does not constrain the far-extreme
($\pm 5$--$7\%$) asymmetry documented above, which an hourly panel cannot resolve.

What survives is worth having. Liquidations predict ETH
tail risk in both tails, symmetrically. Dual permutation nulls (a
circular shift of the shock and an innovation reshuffle) confirm that the
long-horizon coefficient reflects signal rather than overlap noise,
with no signed bias. And the predictive content holds out of sample. In a
rolling-origin pinball-loss evaluation \citep{diebold1995}, adding
liquidation information to an otherwise identical quantile model improves
tail forecasts at both moderate tails at short horizons
(by \OosSkillFiveHone\% and \OosSkillNinetyFiveHone\% at one hour,
$p < 0.01$), though it does not outperform a dedicated GARCH(1,1) scale
model; the claim is incremental information, not dominance. The
volatility channel itself confirms and extends the transaction-level evidence
of \citet{oecd2023} and \citet{lehar2022} to the conditional tails.

The contribution is a portable \emph{methodological caution}, and it
materializes in a new empirical regularity. In quantile local
projections on overlapping cumulative returns, a symmetric volatility
response masquerades as downside fragility. The conjunction of extreme
quantiles, cumulative overlap, and heteroskedastic returns is where the
artifact arises, and the diagnostic battery we provide is built to detect
it. \citet{carriero2024} documented the mechanism for macroeconomic tail
risk; we extend it to asset returns, where the policy narrative it can
mislead (``DeFi liquidations amplify crashes'') is live.

We mean the battery less as a one-off diagnosis of the DeFi data than as a
portable test-bench: a symmetric placebo for shape, dual permutation
nulls for level, a per-period-versus-cumulative exceedance comparison,
standardized skew and equivalence bounds, a cross-asset outcome placebo, and
an out-of-sample evaluation. Any quantile-local-projection study on
overlapping cumulative returns can rerun it to separate scale from shape. That
the apparatus is worth porting is itself a finding. It survives every
adversarial repair we put to it, and its positive control confirms the
equivalence test recovers an asymmetry of actionable size while
staying silent below it, so a bounded null it returns is power-confirmed,
not power-absent.

The
regularity the caution leaves standing is the positive finding: observable
forced deleveraging is a \emph{symmetric} out-of-sample predictor
of ETH tail risk (incremental pinball skill of \OosSkillFiveHone\% and
\OosSkillNinetyFiveHone\% at one hour, $p < 0.01$), to our knowledge among the first quantile-resolution evidence on on-chain forced deleveraging. The
null it delivers is itself informative because it is \emph{bounded}: net of
volatility we exclude any downside-specific amplification larger than
\MdeEightyExc{} per unit of shock at the one-percent tail, a falsifiable statement rather than a
bare failure to reject. The whole is made reproducible by an hourly
on-chain liquidation panel (2021--2025) and a one-command replication
package.

Why does observable forced deleveraging fail to amplify the downside? Scale
and clearing. Median daily DeFi liquidations are \SizeMedianShare{} of one percent of ETH
spot-plus-perpetual turnover; the worst liquidation day of the sample reaches
\SizeShareWorstDay\%; and \citet{moallemi2024} document that over 70\% of
liquidations occur without any downward price jump, consistent with the
near-instant clearing documented by \citet{perez2021} and with the small,
transient pool-level price impacts in \citet{tian2025}. The null is therefore explicitly \emph{size-contingent}.
It characterizes the current scale of observable DeFi deleveraging, not a
structural impossibility, and would bear re-examination if DeFi's share of
crypto leverage grows materially \citep{fsb2023}. Likewise, nothing here
speaks to liquidation cascades on centralized derivatives venues, which our
design cannot observe. The asymmetry of ETH returns is real; our evidence
says only that observable DeFi liquidations show no downside-specific
amplification of it net of volatility.

The paper proceeds as follows. Section~\ref{sec:data} describes the data and
the construction of the liquidation measure. Section~\ref{sec:stylized}
establishes the stylized facts. Section~\ref{sec:method} lays out the locked
specification and inference protocol. Section~\ref{sec:apparent} presents the
naive result. Section~\ref{sec:deconstruction} deconstructs it.
Section~\ref{sec:survives} presents what survives, including the equivalence
bounds and the out-of-sample evaluation. Section~\ref{sec:mechanisms}
discusses mechanisms, Section~\ref{sec:discussion} scope and limitations, and
Section~\ref{sec:conclusion} concludes.